\chapter{总结与展望}

\section{全文总结}

编码边缘计算通过引入冗余缓解了分布式计算中的迟滞效应，但现有研究大多忽略了下行传输阶段由于信道衰落可能导致的通信时延问题。事实上，编码计算在任务分发阶段引入的冗余使得各边缘节点的计算结果之间存在强相关性，这为利用DSC技术优化下行传输提供了可能。本文从DSC的视角出发，研究了面向完整恢复和线性恢复两种场景下的编码计算结果传输方案，主要工作和贡献总结如下：

\begin{enumerate}
	\item[（1）]面向完整恢复的DSC辅助编码计算结果传输方法研究。针对需要完整恢复各子任务计算结果的场景，提出了基于LDPC码伴随式穿刺的DSC方案。该方案利用编码计算结果间的线性相关性，使边缘节点能够根据信道状态自适应调整压缩率，实现传输负载在节点间的灵活重分配。此外，将该方案扩展至高维数据场景，利用逐级编码思想处理多维度数据的压缩与传输。建立了下行传输时延最小化优化模型，设计了基于CCCP的联合优化算法。仿真结果表明，所提方案在信道条件异构场景下显著降低了传输时延，且对迟滞节点数量变化具有良好的鲁棒性。

	\item[（2）]面向线性恢复的DSC辅助编码计算结果传输方法研究。针对仅需恢复计算结果线性函数值的场景，提出了基于嵌套线性码的编码方案。证明了面向线性恢复的编码方案能够获得比经典SW编码更大的可达速率区域，从而实现更高的压缩效率。给出了可达速率区域的简化形式，降低了优化问题的求解复杂度。建立了相应的时延最小化优化模型，设计了联合优化算法。仿真结果验证了所提方案的下行传输性能均优于独立编码方案和SW编码方案，且随着迟滞节点数量增加，所提方案相比SW编码方案的优势更为明显。
\end{enumerate}


\section{未来展望}

本文针对编码计算中可能由于信道衰落导致的通信时延展开了研究，提出了面向完整恢复和面向线性恢复的DSC辅助编码计算结果下行传输方案。虽然提出的方案取得了一定的性能增益，但仍然存在一些问题有待进一步探索。
具体而言，未来可能的探索方向包括：

\begin{enumerate} [ label = （\arabic*）]
    \item 向具有更多边缘节点的密集边缘网络场景扩展。本文在仿真实验中验证了所提方案在中等规模边缘网络中的有效性，然而随着边缘计算技术的发展，未来边缘网络将呈现高密度部署态势。在密集边缘网络场景下，参与计算的边缘节点数量显著增加，这将为所提方案带来新的挑战。首先，大规模节点场景下优化问题的求解复杂度将显著上升，需要设计低复杂度的分布式或启发式求解算法。其次，节点数量的增加意味着更强的用户间干扰，需要研究更先进的多址接入技术或干扰管理策略。

    \item 设计具有更大信源编码可行域的编码方案。本文第三章提出的基于伴随式穿刺的DSC方案仅实现了理想SW区域的一个子区域，如何设计更逼近SW理论极限的实用化编码方案是一个值得探索的方向。一方面，可以研究基于极化码或空间耦合LDPC码的DSC方案，这类码字在有限码长下具有更好的性能逼近特性。另一方面，第四章提出的嵌套线性码方案虽然在特定条件下能够获得比SW区域更大的可达速率区域，但如何扩大该优势的适用条件范围仍有待研究。此外，可以探索机器学习辅助的信源编码方案，利用神经网络学习信源间的相关性结构，实现自适应的压缩编码设计。

    \item 上下行传输链路和计算阶段的联合优化。本文主要聚焦于编码计算的下行结果回传阶段，将上行任务分发和计算阶段视为已知条件。然而，在实际系统中，上行传输、边缘计算和下行传输三个阶段相互耦合、相互影响，联合优化有望进一步提升系统整体性能。具体而言，上行阶段的任务编码方案与下行阶段的信源压缩效率密切相关，不同的编码策略会导致计算结果间相关性的差异，进而影响DSC的压缩增益。此外，计算阶段的资源分配（如CPU频率调度）与通信阶段的功率控制存在耦合关系，需要建立跨计算与通信的联合资源优化框架。因此，研究涵盖上行-计算-下行全链路的联合优化机制具有重要的理论价值和应用前景。
\end{enumerate}
